% Generated draft: requires mathematical and source review.
\documentclass[11pt]{article}
\usepackage[margin=25mm]{geometry}
\usepackage{fontspec}
\setmainfont{Latin Modern Roman}
\usepackage{amsmath,amssymb,mathtools,unicode-math}
\setmathfont{Latin Modern Math}
\usepackage{enumitem,needspace,xurl,hyperref}
\setlength{\parindent}{0pt}
\setlength{\parskip}{5pt}
\setlength{\emergencystretch}{3em}
\newcommand{\N}{\mathbb{N}}
\newcommand{\Z}{\mathbb{Z}}
\newcommand{\Q}{\mathbb{Q}}
\newcommand{\R}{\mathbb{R}}
\newcommand{\C}{\mathbb{C}}
\newcommand{\F}{\mathbb{F}}
\newcommand{\A}{\mathbb{A}}
\newcommand{\PP}{\mathbb{P}}
\newcommand{\E}{\mathbb{E}}
\newcommand{\eps}{\varepsilon}
\newcommand{\dd}{\,\mathrm{d}}
\begin{document}
\section*{UnsolvedMath ID: 30005329}
\textbf{Source problem code:} catalogue entry
\subsection{Definitions and mathematical statement}
The supplied record names a family of type-$B$ $q$-Stirling polynomials $S_q^B[n,k]$, with polynomial variable $q$, but supplies neither their definition nor their admissible index range. It also does not define the property called \emph{parity log-concavity}.

The source assertion, identified as Conjecture 4 and internally attributed to [3, Conjecture 7.5], is:
\begin{quote}
For every admissible pair of indices $(n,k)$, the type-$B$ $q$-Stirling polynomial $S_q^B[n,k]$ is parity log-concave.
\end{quote}
Here both the admissible indices and the meaning of parity log-concavity must be taken from the cited source. The supplied excerpt does not specify the defining polynomials, coefficient conventions, parity subsequences or inequalities, or boundary conventions. Consequently, this is a source-dependent conjectural statement; a self-contained mathematical formulation cannot be recovered from the supplied evidence alone.
\subsection{Short English statement}
The conjecture asks whether every type-$B$ $q$-Stirling polynomial has the source's parity log-concavity property. The supplied record omits the definitions and index conventions needed to express that property precisely.
\subsection{Sources}
\begin{itemize}
\item Supplied local UnsolvedMath registry record, permanent ID $30005329$, source problem code \texttt{catalogue entry}, titled \emph{Parity Log-Concavity of Type-B q-Stirling Polynomials}.
\item \emph{Enumerative Combinatorics (2023)}, Oberwolfach Reports, pp.~3215--3305, DOI \texttt{10.4171/owr/2022/57}: \url{https://doi.org/10.4171/owr/2022/57}. The supplied excerpt labels the assertion Conjecture 4 and attributes it internally to [3, Conjecture 7.5]; the bibliographic identity of reference [3] is not supplied.
\item Supplied literature source: \emph{q-Stirling numbers in type B}, \url{https://arxiv.org/abs/2205.14078}. The record does not establish explicitly whether this is its internal reference [3].
\item Literature assessment dated exactly \texttt{2026-08-22}: the supplied curation labels the conjecture open and reports that it remained explicitly unresolved in the primary treatment located by that cutoff. This is supplied dated curation, not an independently verified literature assessment or verification of the cited claims.
\item The registry records proposed year $2022$ and gives the report citation with year $2023$. No proposer is identified. Its \texttt{corrected\_verified} statement label describes supplied editorial verification of notation and wording; it does not establish mathematical correctness.
\end{itemize}
\end{document}
